\section{Paths as an inductive type indexed by source and target}\label{sec:11-path-algebras:04-paths-as-inductive-type:1}

The key idea for formalizing ``paths compose head-to-tail'' is to make \texttt{Path} an inductive type \emph{indexed} by its own source and target vertex. Lean tracks, in the type itself, which vertex a path starts and ends at. This is the dependent-types idea previewed all the way back in

types indexed by pairs of vertices, exactly the categorical \(\mathrm{Hom}\)-set family \(\mathrm{Hom}_{\mathrm{Free}(Q)}(u,w)\) of the free category on \(Q\).

\begin{lstlisting}
inductive Path {V A : Type} (Q : Quiver V A) : V → V → Type where
  | nil (v : V) : Path Q v v
  | cons {u v w : V} (a : A) (h : Q.source a = v) (h' : Q.target a = w)
      (p : Path Q u v) : Path Q u w
\end{lstlisting}

Reading this carefully:

\begin{itemize}
\tightlist
\item
  \texttt{Path Q : V → V → Type} --- for each pair of vertices \texttt{(start, finish)}, \texttt{Path Q start finish} is the \emph{type} of paths from \texttt{start} to \texttt{finish}. A term of this type is one specific path.
\item
  \texttt{nil (v : V) : Path Q v v} --- the trivial path at \texttt{v}, going from \texttt{v} to \texttt{v} in zero steps. This is \(e_v\) from the math above.
\item
  \texttt{cons \textbackslash{}\{u v w : V\textbackslash{}\} (a : A) (h : Q.source a = v) (h' : Q.target a = w) (p : Path Q u v) : Path Q u w} --- given an existing path \texttt{p} from \texttt{u} to \texttt{v}, and an arrow \texttt{a} whose source is (provably, via \texttt{h}) \texttt{v} and whose target is (provably, via \texttt{h'}) \texttt{w}, \texttt{p} can be extended by appending \texttt{a}, producing a path from \texttt{u} to \texttt{w}. The two proof fields \texttt{h} and \texttt{h'} are exactly what \emph{prevents} the construction of nonsense paths: an arrow cannot be \texttt{cons}ed onto a path unless its source matches where the path left off.
\end{itemize}

Note that \texttt{cons} builds the path by appending the new arrow at the \emph{end}, matching the natural description ``take path \texttt{p}, then follow arrow \texttt{a}.''

\begin{mathreading}

\texttt{Path Q} is the family of Hom-sets of the \textbf{free category} \(\mathrm{Free}(Q)\) on the quiver \(Q\): \texttt{Path Q u w} \(=
\mathrm{Hom}_{\mathrm{Free}(Q)}(u, w)\), the set of directed paths \(u
\rightsquigarrow w\). As an \((V\times V)\)-indexed family it is the map \(V \times V \to \mathbf{Set}\), \((u,w) \mapsto \{\text{paths } u \to w\}\). The constructors are the category structure: \texttt{nil v} is the identity morphism \(e_v = \mathrm{id}_v \in \mathrm{Hom}(v,v)\) (the length-\(0\) path), and \texttt{cons} extends a path by one arrow, with the proof fields \(h : s(a) = v\), \(h' : t(a) = w\) enforcing head-to-tail composability. This is the type-level encoding of ``\(\circ\) is only defined on composable pairs.'' Because the composability constraint lives \emph{in the type}, an ill-formed composite simply fails to typecheck. This is exactly the dependent-type payoff from Chapter 1.

\end{mathreading}

\subsection{\texorpdfstring{Building the path \(\beta\alpha\) (first \(\alpha\), then \(\beta\)) in our example}{Building the path \textbackslash beta\textbackslash alpha (first \textbackslash alpha, then \textbackslash beta) in our example}}\label{sec:11-path-algebras:04-paths-as-inductive-type:2}

\begin{lstlisting}
def pathAlpha : Path exampleQuiver 0 1 :=
  Path.cons ExampleArrow.alpha rfl rfl (Path.nil 0)
\end{lstlisting}

\begin{itemize}
\tightlist
\item
  \texttt{pathAlpha} starts from \texttt{Path.nil 0} (the trivial path at vertex \texttt{0}) and appends \texttt{alpha} (source \texttt{0}, target \texttt{1}, both proved by \texttt{rfl} since they compute directly from the \texttt{match} in \texttt{exampleQuiver}'s definition).
\end{itemize}

\begin{lstlisting}
def pathBetaAlpha : Path exampleQuiver 0 2 :=
  Path.cons ExampleArrow.beta rfl rfl pathAlpha
\end{lstlisting}

\begin{itemize}
\tightlist
\item
  \texttt{pathBetaAlpha} appends \texttt{beta} (source \texttt{1}, target \texttt{2}) onto \texttt{pathAlpha}, giving a path from \texttt{0} to \texttt{2}, exactly \(\beta\alpha\).
\end{itemize}

Attempting \texttt{Path.cons ExampleArrow.beta rfl rfl (Path.nil 0)} instead (appending \texttt{beta}, whose source is \texttt{1}, directly onto the trivial path at \texttt{0}) is rejected by Lean: the proof obligation \texttt{Q.source ExampleArrow.beta = 0} is \texttt{1 = 0}, which is false, so no \texttt{rfl} (or any other proof) exists. This is the type system enforcing ``arrows that do not match up cannot compose'', for free, at compile time.

\begin{mathreading}

\texttt{pathBetaAlpha} is the composite morphism \(\beta\circ\alpha \in \mathrm{Hom}_{\mathrm{Free}(Q)}(0,2)\), built as \(e_0\) followed by \(\alpha : 0\to 1\) followed by \(\beta : 1\to 2\). The \texttt{rfl} proofs discharge the composability side-conditions \(s(\alpha)=0\), \(t(\alpha)=1\), \(s(\beta)=1\), \(t(\beta)=2\), which hold definitionally. The rejected term \(\beta\circ e_0\) fails because it would require \(s(\beta) = 0\), i.e.~\(1 = 0\), a false equation with no proof. This is the type system refusing to compose non-composable arrows, which in a category is the statement that \(\circ\) is a \emph{partial} operation, defined only when endpoints agree.

\end{mathreading}

\textbf{Mathlib equivalent.} Mathlib's \href{https://loogle.lean-lang.org/?q=Quiver.Path}{\texttt{Quiver.Path}} (building on the \texttt{MyArrow}/\texttt{Quiver (Fin 3)} instance from the previous section) is the same inductive family, \texttt{nil}/\texttt{cons}, just with \texttt{cons} taking the shorter path \emph{first} and the new arrow second. This is the mirror image of the book's argument order, which takes the arrow first and the path last:

\begin{lstlisting}
-- NOTE: deliberately *not* `open Quiver` — Mathlib also has a root-level
-- `Path` (a continuous path between two points of a topological space,
-- `Mathlib.Topology.Path`), so an opened `Quiver.Path` would clash with
-- it. Spelling out `Quiver.Path` throughout avoids that.
def pathAlpha' : Quiver.Path (0 : Fin 3) 1 := Quiver.Path.cons Quiver.Path.nil MyArrow.alpha
def pathBetaAlpha' : Quiver.Path (0 : Fin 3) 2 := Quiver.Path.cons pathAlpha' MyArrow.beta
\end{lstlisting}

There is no \texttt{h : Q.source a = v}/\texttt{h' : Q.target a = w} pair to prove here. \texttt{MyArrow.alpha : MyArrow 0 1} (i.e.~\texttt{(0 : Fin 3) ⟶ 1}) already forces the endpoints via its type, exactly as flagged in the previous section. So \texttt{Path.cons} only ever needs the arrow itself, not separate proofs about it.
